% !TeX root = ../main.tex

\section{Conclusions}

The finite difference method combined with sparse matrices to perform the diagonalization results in a powerful technique for solving the 2D Schrödinger equation in arbitrary geometries. Using this approach, we obtain the eigenvalue spacings and observe signatures of chaotic behavior. \\

We show that the level spacing distribution in the rectangular billiard is Poissonian, which can be interpreted as the eigenvalues being uncorrelated. As a consequence, there is a high probability of finding two eigenvalues very close to each other or even degenerate. On the other hand, the spacing distribution in the stadium billiard follows Wigner-Dyson statistics, which indicates that the eigenvalues are correlated. Therefore, the probability of finding two very close or degenerate eigenvalues is strongly suppressed. \\

It is important to note that the complete stadium billiard possesses spatial symmetries (even and odd parity), causing the spectrum to become a superposition of independent subsystems. By considering the quarter-stadium geometry, these symmetry sectors are separated, allowing the pure Wigner-Dyson statistics to emerge without contamination from states of different parity. This isolation of the spatial symmetries is essential for obtaining undistorted spectral statistics that properly reflect the chaotic nature of the system. \\

In summary, this work successfully demonstrates how the complex statistical properties predicted by Random Matrix Theory naturally emerge when numerically solving a simple physical system under specific geometric boundary conditions.

